% =============================================================================
% Section 12 — Open-Source Strategy
% Author: David (Research Lead)
% =============================================================================

\section{Open-Source Strategy}
\label{sec:oss-strategy}

\subsection{Is This a Viable Open-Source Project?}

The short answer is yes, but with conditions.
The diagram-as-code category has demonstrated commercial and community
viability at scale: Mermaid~\cite{mermaid2023} has over 88{,}000 GitHub
stars, raised \$7.5M in seed funding in 2024~\cite{mermaidchart2024}, and
has been natively integrated into GitHub, GitLab, Notion, Obsidian, and
Azure DevOps.
The lesson from Mermaid is that a plain-text format with zero-install
rendering, embedded in documentation ecosystems, can achieve massive adoption.

Timeline Compiler's gap analysis (Section~\ref{sec:comparison}) identifies
a cell in the tool landscape that is genuinely unoccupied: an
\emph{open-source}, \emph{presentation-quality}, \emph{agent-generation-friendly}
timeline compiler.
This gap is both an opportunity and a risk; the following sections analyse
both.

% -----------------------------------------------------------------------------
\subsection{Closest Existing Projects and the Gap They Leave}
\label{sec:oss-gap}
% -----------------------------------------------------------------------------

\begin{itemize}
  \item \textbf{Mermaid} (\S\ref{sec:comparison}) leaves the presentation-quality
        gap.  Its \texttt{timeline} and \texttt{gantt} types are documentation
        tools, not slide-production tools.  No PPTX target.  No theme engine.
        Layout is viewport-dependent, not deterministic.

  \item \textbf{TimelineJS}~\cite{timelinejs} (Knight Lab, GPL) is the
        best-known open-source timeline tool, but it is a web-storytelling
        widget (iframes, Google Sheets), not a document/slide compiler.
        No SVG/PPTX/PDF output; no CLI; no agent path.

  \item \textbf{Frappe Gantt}~\cite{frappegantt} and
        \textbf{vis-timeline}~\cite{vistimeline} are interactive browser
        widgets.  No declarative text format; no static artefact output.

  \item \textbf{think-cell}~\cite{thinkcell} produces the closest visual
        quality but is closed-source, expensive, and PowerPoint-locked.
        It is the benchmark for what Timeline Compiler output should look like,
        not a substitute for it.

  \item \textbf{Plotly.js}~\cite{plotlyjs} can produce high-quality timeline
        charts, but requires programming and has no declarative timeline
        grammar or agent-generation story.
\end{itemize}

\noindent
The gap is real: the market has diagram-as-code (low quality, wrong idiom) or
presentation tools (high quality, closed, manual).
Nothing sits in between.

% -----------------------------------------------------------------------------
\subsection{Potential Adopters}
\label{sec:oss-adopters}
% -----------------------------------------------------------------------------

Four adopter personas have been identified based on real-data patterns observed
during the research phase:

\paragraph{Developer-Relations and Technical Writers.}
Dev-rel teams already live in Mermaid and Markdown.
They need to produce roadmaps and architecture-evolution timelines for
conference talks, blog posts, and product docs.
The barrier to adoption is near-zero if Timeline Compiler has a CLI, a
Markdown code-fence syntax, and SVG output.
This group provides early adopters and community feedback.

\paragraph{Product Managers and Programme Leads.}
PMs need quarterly roadmaps and programme timelines for steering committees
and executive reviews.
They currently use PowerPoint or think-cell manually.
A compiler that takes a YAML/CSV definition and emits a PPTX with
think-cell-quality layout would replace hours of work per roadmap cycle.
This group provides the strongest word-of-mouth growth vector.

\paragraph{Management Consultants (Freelance / Boutique).}
Large firms use think-cell.
Freelancers and boutique consultancies cannot justify \$27.30/user/month
per project.
An open-source tool with a polished PPTX target has a credible value
proposition here.

\paragraph{AI Agent Toolchains.}
The emergent user is not a human at all.
LLM-based planning agents (GitHub Copilot, OpenAI Assistants, Anthropic
Claude with tools) generate project plans, architecture proposals, and
transformation roadmaps.
Today there is no standard open-source \emph{target format} for those plans.
Timeline Compiler's IR, if published as a JSON Schema and exposed as an MCP
tool~\cite{mcp-spec}, becomes the target format AI agents compile to.
This persona drives long-tail adoption as agent toolchains proliferate.

% -----------------------------------------------------------------------------
\subsection{Ecosystem Fit}
\label{sec:oss-ecosystem}
% -----------------------------------------------------------------------------

\subsubsection{Mermaid / Markdown / CI Ecosystem}

Timeline Compiler should position itself as the ``next step'' after Mermaid:
use Mermaid in your README; use Timeline Compiler when you need a slide.
Concretely:

\begin{itemize}
  \item A Mermaid-syntax ingester (accepting \texttt{timeline} and
        \texttt{gantt} blocks) provides a zero-friction migration path.
  \item A Markdown code-fence processor (similar to Mermaid's own fence
        renderer) enables embedding timelines in MkDocs, VitePress, and
        Quarto~\cite{mermaid2023}.
  \item A GitHub Actions step for CI-driven rendering integrates into
        the workflow teams already use for diagram-as-code.
  \item SVG output embeds directly in Slidev~\cite{slidev} and
        Obsidian~\cite{obsidian} --- both are Mermaid-native environments.
\end{itemize}

\subsubsection{MCP / Agent Ecosystem}

The Model Context Protocol~\cite{mcp-spec} is the emerging standard for
LLM tool invocation.
An MCP server that exposes Timeline Compiler as a tool would allow any
MCP-compatible agent (GitHub Copilot, Claude, etc.) to:

\begin{enumerate}
  \item Accept a natural-language plan from the user.
  \item Generate a valid Timeline IR (JSON/YAML, constrained by the
        published JSON Schema).
  \item Invoke the Timeline Compiler MCP tool.
  \item Receive back an SVG, PNG, or PPTX and embed it in the agent's
        response.
\end{enumerate}

\noindent
OpenAI's Structured Outputs feature~\cite{openai-structured-outputs}
(JSON Schema-constrained generation) enables agents to reliably produce
valid IR without hallucinating schema fields.
The design implication is that the Timeline IR schema \emph{must} be
published as a JSON Schema document alongside the compiler.

\subsubsection{PPTX Ecosystem}

python-pptx~\cite{python-pptx} is the leading open-source Python library for
programmatic PPTX generation.
It is MIT-licensed, well-maintained, and can produce slides that open
natively in PowerPoint and Keynote.
Building the PPTX output target on python-pptx avoids a dependency on
LibreOffice or headless Chrome, making the compiler pip-installable with
minimal system requirements.

% -----------------------------------------------------------------------------
\subsection{Extension Opportunities}
\label{sec:oss-extensions}
% -----------------------------------------------------------------------------

The compiler-IR architecture creates natural extension points:

\begin{description}
  \item[Ingesters.] New data sources can be added without touching the
        renderer: CSV/Excel, Jira API, Linear API, Notion database, ADO
        work items~\cite{ado-workitems}, GitHub Projects~\cite{github-projects},
        markdown plans, Mermaid blocks.
        Ingesters are the highest-traffic contribution path for community
        members.

  \item[Themes.] A typed theme schema (fonts, colour palettes, swimlane
        heights, milestone shapes) allows community contribution of brand
        themes without touching rendering code.
        McKinsey-style, BCG-style, GitHub-style, and corporate-brand themes
        are natural community contributions.

  \item[Output Targets.] Additional renderers (HTML interactivity, Keynote,
        Google Slides API, LaTeX beamer) can be added without changing the
        IR or ingesters.
        The Vega-Lite model~\cite{vegalite2017} demonstrates that a clean
        IR/renderer separation enables this pattern reliably.

  \item[Grammar Extensions.] The IR can be versioned and extended (new event
        types, annotations, dependency arrows) without breaking existing
        themes or renderers, following established IR versioning practice
        from compiler design~\cite{dragonbook2006}.
\end{description}

% -----------------------------------------------------------------------------
\subsection{Strongest Differentiators}
\label{sec:oss-differentiators}
% -----------------------------------------------------------------------------

Three differentiators are assessed as strongest, in decreasing order of
defensibility:

\begin{enumerate}
  \item \textbf{The only open-source, agent-generation-friendly timeline
        compiler with PPTX output.}
        No existing open-source tool combines a stable text-format IR,
        deterministic rendering, and a PPTX output target.
        This is the clearest gap in the landscape.

  \item \textbf{Visual-communication grammar vs.\ task-scheduling grammar.}
        Positioning the tool as a ``timeline grammar for human communication''
        rather than a ``project management tool'' is a genuine conceptual
        distinction.
        It attracts users who are currently underserved by Mermaid
        (insufficient quality) and repelled by MS Project (wrong abstraction).

  \item \textbf{AI-agent-first design.}
        Publishing a JSON Schema for the IR and an MCP tool interface makes
        Timeline Compiler the natural compiler target for any LLM that
        generates plans or roadmaps.
        As agent-driven workflows proliferate~\cite{mcp-spec}, this
        positions Timeline Compiler at the centre of a growing
        distribution channel with no human adoption friction.
\end{enumerate}

% -----------------------------------------------------------------------------
\subsection{Community Contribution Model}
\label{sec:oss-community}
% -----------------------------------------------------------------------------

The Mermaid model~\cite{mermaidchart2024} is instructive: open-source core
(MIT), commercial SaaS layer for hosted rendering, team features, and
enterprise integrations.
Timeline Compiler should follow a similar open-core pattern:

\begin{itemize}
  \item \textbf{Core compiler} (MIT): IR schema, renderer, CLI, standard
        output targets, standard ingesters.
        All of this is open-source.
        Community contributions to ingesters and themes are accepted here.

  \item \textbf{Theme marketplace} (community-driven):
        A GitHub-hosted registry of community themes, analogous to the
        npm registry for packages or the VS Code extension marketplace.
        Themes are independent packages; they do not require merging into
        the core.

  \item \textbf{Hosted service} (future, optional):
        Web UI for non-technical users, CI integration hooks,
        and enterprise SSO.
        This funds maintenance without restricting the open-source core.
\end{itemize}

\noindent
The contribution surface most likely to attract community contributions
(in descending order of probability) is: ingesters, then themes, then
additional output targets, then IR schema extensions.
Core renderer changes require higher trust and should require maintainer
review.

% -----------------------------------------------------------------------------
\subsection{Adoption Risk Assessment}
\label{sec:oss-risk}
% -----------------------------------------------------------------------------

No adoption analysis is honest without a candid risk section.

\paragraph{Risk 1: The presentation-quality bar is high.}
think-cell is the quality benchmark.
Matching it with an open-source, text-driven renderer is a significant
engineering challenge.
If the initial PPTX output is visually inferior to think-cell, the
``presentation quality'' differentiator collapses, and Timeline Compiler
is merely a less capable Mermaid.
\textit{Mitigation}: Define a strict visual quality bar in the design
spec and validate against real executive roadmaps before launch.

\paragraph{Risk 2: The IR may be too narrow or too wide.}
An IR that is too narrow (only covering the simplest roadmaps) will be
inadequate for real-world use.
An IR that is too wide (attempting to model every possible timeline idiom)
will be complex to understand, hard for agents to generate reliably, and
fragile to maintain.
\textit{Mitigation}: Scope the IR to the five canonical timeline types
identified in the scope section, and use real data crawls (ADO exports,
GitHub Projects exports) to validate coverage before freezing the schema.

\paragraph{Risk 3: Mermaid is a gravity well.}
Mermaid has 88{,}000 stars, is embedded everywhere, and is improving.
If Mermaid's \texttt{timeline} type adds PPTX export and better theming,
the core differentiator weakens.
\textit{Mitigation}: The visual-communication grammar thesis and the
agent-generation design are harder to retrofit into Mermaid than a PPTX
renderer. The IR schema and MCP interface are the durable moats.

\paragraph{Risk 4: Agent toolchains may standardise on Mermaid instead.}
LLMs already reliably generate Mermaid syntax.
If the agent ecosystem converges on Mermaid + post-processing as the
standard agent-to-timeline pipeline, the agent-first positioning loses force.
\textit{Mitigation}: Publish the JSON Schema early; build the MCP tool
before the ecosystem settles.  Structured Outputs~\cite{openai-structured-outputs}
are demonstrably more reliable than free-text Mermaid generation for
complex schemas.

\paragraph{Risk 5: PPTX format instability.}
The OOXML PPTX specification is controlled by Microsoft.
python-pptx~\cite{python-pptx} may lag new PowerPoint features.
\textit{Mitigation}: Treat PPTX as an output layer, not the IR.
SVG and PDF are canonical; PPTX is a convenience export.
Degrade gracefully when PPTX features are unavailable.
